% Round 18. Author: general_u2 branch. New interface, not a new Green family.
\section{Primitive representatives of the full numerical $U_2$ defect}
\label{r18pu:section}

\paragraph{Historical scope before stating the interface.}
The divisor-kernel formula, ample-weighted injectivity on its image,
and primitive normalization of an \emph{actual generic homological
cycle} already occur in the frozen R7 sources.  So do the total
isotropy and rational Witt-index bound for the actual obstruction
image.  They are not claimed as new results of this round.
The purpose here is to identify the \emph{entire model numerical
space} $U_2$ introduced in DQ17 with a precisely specified primitive
subspace, and then to identify the actual period map with its
possibly proper cycle-generated subspace.  This is a current-interface
identification and an independently checked structural strengthening,
not a new class of generic fibers or a solution of general G.

The exact unchanged historical sources, with their archive paths
and SHA-256 hashes, are recorded in
\path{u2/pu18_historical_bridge.json}.  They are:
\begin{itemize}
\item R7 \path{threefold_elementary_kernel.tex}, Theorem
      \texttt{bb7tf:elementary-kernel}, equation TF.E1;
\item R7 \path{threefold_codimension_two_obstruction.tex},
      Proposition \texttt{bb7tf:weighted-kernel}, equation TF.P1,
      and the subsequent actual-cycle normalization;
\item R7 \path{threefold_boundary_primitive_obstruction.tex},
      TF.O1--TF.O6, concerning the actual obstruction image;
\item R7 \path{descent_inertia_primitive_remainder.tex},
      R7.IR9--R7.IR15, concerning a separate cohomological quotient
      under the displayed divisor-Tate hypothesis.
\end{itemize}
The last cohomological quotient is not identified with $U_2$ here.
The old draft headers in those unchanged files do not alter the
later frozen registers.

\subsection{A complete check of the divisor input}

Use the reduced strict semistable native model of DQ17, of relative
dimension $n\ge3$.  Fix one relatively ample rational class $H$.
Let $T^0$ be the direct sum of the fundamental-class lines of the
native connected triple strata, and let
\[
 \delta_0:W^0\longrightarrow T^0,\qquad
       (\delta_0b)_{ijk}=b_{ij}-b_{ik}+b_{jk}.
 \tag{PU18.1}\label{r18pu:delta}
\]
All restrictions to disconnected strata are taken on their actual
native connected pieces.  Thus multiple edges and multiple triples
are included, with their ordinary signed incidence maps.
Put
\[
 \begin{gathered}
 q_W(x,y)=\sum_e\deg_{E_e}(x_ey_eH_e^{n-3}),\qquad
 R_1=\operatorname{im}\rho_1,\\
 d_0=\rho_1\sigma_1:W^0\to W^1,\qquad J=\operatorname{im}d_0 .
 \end{gathered}
 \tag{PU18.2}\label{r18pu:edge-form}
\]
Divisor Hodge index makes $q_W$ nondegenerate, with positive
index equal to the number of native connected edges.

\begin{lemma}[The full radical, including harmonic edge directions]
\label{r18pu:radical}
One has
\[
                       \operatorname{rad}(q_W|_{R_1})=J.
 \tag{PU18.3}\label{r18pu:radical-equality}
\]
Consequently
\[
 \ker(H^{n-3}A_1)=\ker A_1=B_1+L_1,
 \qquad
 H^{n-3}:\operatorname{im}A_1\longrightarrow C^{n-1}
                           \text{ is injective}.
 \tag{PU18.4}\label{r18pu:weighted-kernel}
\]
These conclusions reproduce the R7 divisor input in current signs.
\end{lemma}

\begin{proof}
The degree-zero instance of DG18's actual triple factorization gives
$d_0=\tau_1\delta_0$, where $\tau_1$ is the alternating triple Gysin
into edge divisors.  Projection formula therefore gives
\[
 q_W(Hb,d_0c)=
      \sum_T\deg_T(H_T^{n-2})(\delta_0b)_T(\delta_0c)_T .
 \tag{PU18.5}\label{r18pu:scalar-energy}
\]
Every coefficient in this scalar form is positive.  It follows,
by also setting $c=b$, that
\[
       \ker d_0=\ker\delta_0,\qquad
       \dim J=\operatorname{rank}\delta_0,\qquad
       HW^0\cap J^\perp=H\ker\delta_0 .
 \tag{PU18.6}\label{r18pu:J}
\]
Here the last perpendicular is for $q_W$.  On $HW^0$ the form
is positive definite, with diagonal entries
$\deg_E(H_E^{n-1})$.

DQ17's identity $A_1\sigma_1=0$ and its energy adjunction give
$q_W(d_0b,\rho_1v)=0$ for every $b,v$.
Thus $J\subset R_1\cap R_1^\perp$; in particular $J$ is totally
isotropic.

DG18's unconditional $U_1=0$ shows
\[
                 \ker\delta_0=\operatorname{im}\rho_0
                                      \oplus\ker\sigma_1 .
 \tag{PU18.7}\label{r18pu:harmonic-split}
\]
Indeed, for $b\in\ker d_0$, the class $\sigma_1b$ lies in
$B_1\cap L_1=\operatorname{im}A_0$.  Write it as
$\sigma_1\rho_0c$; then $b-\rho_0c\in\ker\sigma_1$.
If $\rho_0c\in\ker\sigma_1$, then $A_0c=0$.
The native weighted graph-Laplacian energy, whose edge weights
are $\deg_E(H_E^{n-1})>0$, gives $\rho_0c=0$.  This proves
the directness of the sum.  No vanishing of the first cohomology
of the dual incidence complex has been assumed.

The two positive subspaces
$H\operatorname{im}\rho_0$ and $H\ker\sigma_1$ in
\eqref{r18pu:harmonic-split} are orthogonal.  Projection formula
proves this, since pairing $Hb$ with $H\rho_0c$ is the degree
pairing of $\sigma_1b$ against the corresponding component
multiple of $H^{n-1}$.  More generally
\[
                    H\ker\sigma_1\subset R_1^\perp,
             \qquad H\operatorname{im}\rho_0\subset R_1.
 \tag{PU18.8}\label{r18pu:two-subspaces}
\]

Pass to the nondegenerate real quadratic space $J^\perp/J$.
Its positive index is
$\#\{\text{native edges}\}-\dim J=\dim\ker\delta_0$.
The image of the positive definite space $H\ker\delta_0$
is injective and has exactly that dimension, by
\eqref{r18pu:J}.  It is therefore a maximal positive subspace,
and its perpendicular is negative definite.
The subspace $R_1/J$ contains
$H\operatorname{im}\rho_0$ and is orthogonal to
$H\ker\sigma_1$.  Decomposing with respect to the two orthogonal
positive subspaces in \eqref{r18pu:harmonic-split} shows that
$R_1/J$ is the orthogonal direct sum of
$H\operatorname{im}\rho_0$ and a subspace of that negative
definite perpendicular.  Its restricted form is nondegenerate.
This proves \eqref{r18pu:radical-equality}, first over $\mathbf R$
and then over $\mathbf Q$.

For $F=H^{n-3}A_1:C^1\to C^{n-1}$, the native perfect
pairing of these two spaces gives
\[
                 \langle Fv,u\rangle_C
                         =-q_W(\rho_1v,\rho_1u).
 \tag{PU18.9}\label{r18pu:F}
\]
Thus $\ker F=\rho_1^{-1}(J)=B_1+L_1$.
DQ17 already gives $B_1+L_1\subseteq\ker A_1$, whereas
$\ker A_1\subseteq\ker F$.  This proves both equalities and
the stated injection on $\operatorname{im}A_1$.
\end{proof}

\subsection{The entire numerical defect has a unique primitive representative}

Now let $n\ge4$ and set
\[
 P_H^2(C)=\ker(H^{n-3}:C^2\to C^{n-1}),\qquad
 \mathscr U_H=P_H^2(C)\cap B_2\cap L_2 .
 \tag{PU18.10}\label{r18pu:primitive-space}
\]

\begin{theorem}[A rational isomorphism for the full $U_2$]
\label{r18pu:full-defect}
The quotient map induces a canonical rational isomorphism
\[
          \mathscr U_H\xrightarrow{\sim}
             U_2=(B_2\cap L_2)/\operatorname{im}A_1 .
 \tag{PU18.11}\label{r18pu:U-isomorphism}
\]
Equivalently,
\[
                    B_2\cap L_2
                       =\operatorname{im}A_1\oplus\mathscr U_H.
 \tag{PU18.12}\label{r18pu:direct-sum}
\]
The dependence on $H$ is the explicitly specified primitive
representative, while the quotient $U_2$ itself does not depend on it.
\end{theorem}

\begin{proof}
For $z\in B_2\cap L_2$, the functional
$\lambda_z(v)=\langle H^{n-3}z,v\rangle_C$ on $C^1$
annihilates $B_1+L_1$.
For $v\in B_1$, the class $H^{n-3}v$ is in $B_{n-2}$,
which is orthogonal to $L_2$ by signed adjunction.
For $v\in L_1$, its ample multiple is in $L_{n-2}$,
which is orthogonal to $B_2$.  These are two distinct
uses of the two numerator conditions on $z$.

The rational map $F$ in \eqref{r18pu:F} is symmetric, and
Lemma~\ref{r18pu:radical} identifies its kernel with $B_1+L_1$.
Its image is the annihilator of that kernel.  Hence there
is $v\in C^1$ with $Fv=H^{n-3}z$.  Put
$z_H=z-A_1v$.  Since $\operatorname{im}A_1\subseteq B_2\cap L_2$,
this belongs to $\mathscr U_H$ and represents the original
quotient class.

If two such representatives differed, their difference would
be $A_1v$ with $H^{n-3}A_1v=0$.  The last assertion of
\eqref{r18pu:weighted-kernel} makes that difference zero.
This proves existence, uniqueness and the direct sum.
Every matrix and right-hand side is rational; no scalar
extension of cycle coefficients is used.
\end{proof}

\subsection{A bound for the whole fourfold numerical defect}

In relative dimension four set
\[
 q_C(z,z')=\sum_i\deg_{D_i}(z_i z_i'),\qquad
 \rho_i=\dim N^1(D_i)_{\mathbf Q},\qquad
 \eta_i=\dim N^2(D_i)_{\mathbf Q}-\rho_i .
 \tag{PU18.13}\label{r18pu:fourfold-form}
\]
For each component let $b_i^-$ be the negative index of its
nondegenerate middle numerical form, and put
\[
                          \epsilon_i=b_i^--(\rho_i-1).
 \tag{PU18.14}\label{r18pu:excess-negative}
\]

\begin{corollary}[Primitive Witt and signature bounds on $U_2$]
\label{r18pu:Witt}
The form $q_C$ on $P_H^2(C)$ is nondegenerate, of dimension
$\sum_i\eta_i$ and negative index $\sum_i\epsilon_i$.
In particular each $\epsilon_i\ge0$.
The subspace $\mathscr U_H$ is totally isotropic, and
\[
 \begin{split}
 \dim_{\mathbf Q}U_2
  &\le \operatorname{wi}_{\mathbf Q}(q_C|_{P_H^2(C)})\\
  &\le \min\left\{\sum_i\epsilon_i,\
                         \sum_i(\eta_i-\epsilon_i)\right\}.
 \end{split}
 \tag{PU18.15}\label{r18pu:Witt-bound}
\]
Here the rational Witt index means the largest dimension
of a rational totally isotropic subspace.
Thus a nonzero $U_2$ requires a rational isotropic primitive
tuple and requires a failure of primitive middle positivity
on at least one component.  Neither condition is asserted
to be sufficient for $U_2\ne0$.
\end{corollary}

\begin{proof}
Divisor Hodge index makes the form on $HN^1(D_i)$
nondegenerate with signature $(1,\rho_i-1)$.
Its orthogonal complement in $N^2(D_i)$ is precisely
$\ker(H:N^2(D_i)\to N^3(D_i))$, by numerical duality
between $N^3(D_i)$ and $N^1(D_i)$.
This proves the dimension, nondegeneracy and negative-index
statements, including $\epsilon_i\ge0$.

For $z,z'\in\mathscr U_H$, one lies in $L_2$ and the
other in $B_2=L_2^\perp$, so $q_C(z,z')=0$.
The isomorphism \eqref{r18pu:U-isomorphism} therefore
places the entire $U_2$, not merely a cycle-image quotient,
in a rational totally isotropic subspace of this primitive
form.  The first inequality follows by definition of
Witt index; extension to $\mathbf R$ gives the usual bound
by the smaller of the positive and negative indices.
\end{proof}

When every $\epsilon_i=0$, this supplies another proof of the
fourfold numerical $U_2$ consequence of S18, without invoking the
full BGS stratum theorem.  It does not improve the R7 actual-cycle
anisotropy criterion to a new Green domain: that earlier criterion
already applies directly to the actual primitive obstruction
image, without the extra double-comparison assumption needed
for the next identification.

\subsection{The actual cycle image can be strictly smaller}

Assume $n=4$ and $K_2=B_2$, for example by the precise CD17
curve-comparison hypothesis on all doubles.  For a geometrically
homologically trivial native class $z\in\operatorname{CH}^2(X)_{\mathbf Q}$,
choose a Chow extension with its frozen pure operational
representative, and write $c_z\in C^2$ for its numerical
component restriction tuple.

\begin{proposition}[The cycle period map in the primitive coordinates]
\label{r18pu:period-image}
One has $c_z\in B_2\cap L_2$.  Its class modulo
$\operatorname{im}A_1$, and hence its unique primitive
representative $(c_z)_H\in\mathscr U_H$, is independent of the
chosen extension.  The resulting rational linear map
\[
 \operatorname{CH}^2(X)_{\mathbf Q,\mathrm{geom\,hom}}
              \longrightarrow \mathscr U_H,\qquad z\longmapsto(c_z)_H
 \tag{PU18.16}\label{r18pu:cycle-map}
\]
has kernel exactly the strong-G subgroup of VT17.  Its image
is a subspace of $\mathscr U_H$; surjectivity is not asserted.
For $v\in V_2$, its native normalized model period is
\[
                      \nu_z(v)=q_C((c_z)_H,v).
 \tag{PU18.17}\label{r18pu:period-formula}
\]
This is DQ17's perfect duality with
$V_2/(B_2+L_2)$, expressed in primitive coordinates.
\end{proposition}

\begin{proof}
The two further restrictions of one operational model class
to a double agree in actual Chow.  Thus $c_z\in B_2$.
ER17/FD17 kills its period against $K_2=B_2$ because $z$
is homologically trivial.  Under the perfect middle
numerical pairing, $B_2^\perp=L_2$ by DQ17's signed
adjunction.  Therefore $c_z\in L_2$ as well.

Localization changes the component tuple of two extensions
of the same generic Chow class by an element of
$\operatorname{im}A_1$.  Theorem~\ref{r18pu:full-defect}
gives the stated independence and linearity.
By definition the native model period is $q_C(c_z,v)$.
The difference $c_z-(c_z)_H$ lies in $\operatorname{im}A_1$,
which is orthogonal to $V_2=\ker A_2$ by DQ17.
This proves \eqref{r18pu:period-formula}.

All periods on $V_2$ vanish precisely when
$c_z\in V_2^\perp=\operatorname{im}A_1$, precisely when
$(c_z)_H=0$.  FD17 and the necessary-and-sufficient theorem
VT17 identify that zero locus with existence of the
ordinary actual Green form on the original specified
support complement, tested against all original compact
actual forms.  The primitive tuple is an algebraic
description of the obstruction; it does not replace the
analytic output or narrow the test space.
\end{proof}

\paragraph{The exact remaining question.}
For a fourfold with the stated double comparison, general G asks
that the image in \eqref{r18pu:cycle-map} be zero.  Vanishing of
the full $\mathscr U_H$, equivalently $U_2=0$, is stronger.
Even if the full numerical model space were nonzero, its actual
homological cycle image could vanish.  Without double comparison
there is also the independent CL17 compatible-representative
defect; \eqref{r18pu:cycle-map} has not removed it.

This distinction is also why a rational isotropic vector in a
component middle form is not a counterexample to G: it must
first belong to the compatible signed edge image, and then
must be realized by the component restrictions of a generic
homologically trivial cycle.  Neither realization is supplied
by quadratic-form linear algebra alone.

The new model theorem PG18 proves $U_2=0$ in its common-ambient
product-smoothing class by an actual triple calculation.
The present section locates the general remaining numerical
defect precisely, but adds no unrestricted G or height theorem.
